% Build twice with pdfLaTeX. verify.py generates the table rows from saved results.
\documentclass[11pt,letterpaper]{article}
\usepackage[T1]{fontenc}
\usepackage{lmodern,amsmath,amssymb,booktabs,xcolor}
\usepackage[letterpaper,top=.79in,bottom=.78in,left=.98in,right=.98in,footskip=.33in]{geometry}
\usepackage{hyperref}
\definecolor{teal}{HTML}{145F63}
\hypersetup{colorlinks=true,linkcolor=teal,citecolor=teal,urlcolor=teal,
 pdftitle={Self-review of real-analysis proofs},pdfauthor={Vivaan Pinglikar},
 pdfsubject={A local Qwen3.8-27B pilot benchmark with a Lean audit}}
\urlstyle{same}
\setlength{\parindent}{1.15em}
\setlength{\parskip}{2.5pt}
\setlength{\emergencystretch}{1.7em}
\raggedbottom
\newcommand{\cut}{\textsuperscript{*}}
\newcommand{\lead}[1]{\par\smallskip\noindent\textbf{#1}\enspace\ignorespaces}
\newcommand{\code}[1]{\texttt{\detokenize{#1}}}
\newcommand{\quoted}[1]{\textquotedblleft{}#1\textquotedblright{}}
% Generated by verify.py from saved individual results.
\newcommand{\resultrows}{%
First attempt & 6/10 & --- & 8/10 \\
Neutral review & 8/10 & 2/4 (50\%) & 6/10 \\
Flagged review & 0/4 & 0/4 (0\%) & 4/4 \\
Fresh retry & 1/4 & 1/4 (25\%) & 4/4 \\
}
\newcommand{\costrows}{%
First attempt & 10 & 1,523 & 19,744 & 8 \\
Neutral review & 10 & 21,545 & 18,248 & 6 \\
Flagged review & 4 & 8,918 & 8,192 & 4 \\
Fresh retry & 4 & 666 & 8,192 & 4 \\
Total & 28 & 32,652 & 54,376 & 22 \\
}
\newcommand{\itemrows}{%
P01 & Means of a convergent sequence & Valid\cut & Valid & --- & --- \\
P02 & Monotonicity and sparse spikes & Gap\cut & Valid\cut & Gap\cut & Valid\cut \\
P03 & Asymptotic continuity & Valid\cut & Valid\cut & --- & --- \\
P04 & A moving maximum & Valid\cut & Valid & --- & --- \\
P05 & A common intersection & Valid & Valid & --- & --- \\
P06 & A derivative at one point & Valid & Valid & --- & --- \\
P07 & Weighted concentration & Valid\cut & Valid\cut & --- & --- \\
P08 & Moving observation points & Gap\cut & Gap\cut & Gap\cut & Gap\cut \\
P09 & Matching a secant with a derivative & Gap\cut & Gap\cut & Gap\cut & Gap\cut \\
P10 & Dense zeros and the Riemann integral & Gap\cut & Valid\cut & Gap\cut & Gap\cut \\
}

\begin{document}
\begin{center}
  {\LARGE\bfseries Self-review of real-analysis proofs\par}
  \vspace{3pt}
  {\large A local Qwen3.8-27B pilot benchmark with a Lean audit\par}
  \vspace{7pt}
  {\large Vivaan Pinglikar\par}
  {\small\itshape Independent researcher\par}
  {\small November 2026\par}
\end{center}
\begin{abstract}
I tested whether asking Qwen to check a proof helped it correct mistakes.
On ten handpicked problems, a neutral request to review its work raised the
score from 6/10 to 8/10. Of the four first attempts graded incorrect, neutral
review repaired two, a fresh retry repaired one, and review with an explicit
error flag repaired none. A retrospective audit used Lean to check solutions
to all ten problems and retained all 28 prose grades. Every follow-up on an
initial error reached the output limit. The extra repair from neutral review
came from one problem.
\end{abstract}

\section{Benchmark design}
I ran all 28 generations locally on my computer using LM Studio. The ten
problems cover sequences and series, continuity, uniform convergence,
compactness, differentiation and integration. Three were rated easy, four
medium and three hard. They were chosen to span topics and difficulty.
Variants of these classical exercises may have appeared in the model's
training data.

Each first attempt used the same instruction and a fresh context. \emph{Neutral
review} (A) gave Qwen the problem and its first answer and asked it to check each
step. Every problem received A. For the four initial errors, \emph{flagged
review} (B) supplied the same answer and an incorrect verdict without the
grader's reason. A \emph{fresh retry} (C) gave only the original problem and
first-attempt instructions. Each follow-up started in a new context with the
same output budget and sampling settings. The three follow-ups on a given
problem shared a second-round seed. Reference proofs and rubrics were private
during generation.

\section{Results}
\begin{center}
\small\renewcommand{\arraystretch}{1.13}
\begin{tabular}{@{}lrrr@{}}
\toprule
Condition & Correct / attempted & Initial errors repaired & Cutoffs / attempted\\
\midrule
\resultrows
\bottomrule
\end{tabular}
\end{center}

Neutral review preserved all six initially correct answers (0/6 broken).
For the paired comparison of flagged review with a fresh retry, both fixed
zero problems, flagged review alone fixed zero, fresh retry alone fixed one
(P02), and neither fixed three. B versus C changes both access to the first
answer and the wording of the instruction. A versus B on the four errors
holds access to the first answer fixed.

\clearpage
\section{Grading and formal verification}
One AI grader assessed the entire visible response, including mathematics
outside the requested proof tags. Any sound method could earn credit. A
material gap, circular reasoning, an unjustified theorem or an omitted
part earned an incorrect grade. The rubric allowed routine algebra to be
implicit. Second-round responses were shuffled and their condition labels
hidden before grading. Some wording could still suggest a condition.

In a separate retrospective audit, Lean~4.19.0 checked formal solutions to
all ten questions, including the requested examples. The project contains
37 named theorems, with their assumptions and machine-checkable proofs.
The AI audit compared the archived English answers with those arguments and
the original rubric. It had access to the earlier grades and condition labels.
Lean verifies that a formal argument follows from its assumptions. Deciding
whether a sentence in Qwen's answer supplied a necessary step still required
judgment. The audit retained all 28 grades. Independent human grading remains
outstanding.

\lead{P10. Why neutral review earned the extra repair.}
A graph can reach height one at a single point and still have zero area. The
problem also asked why \emph{every} nonnegative, Riemann integrable function
with a zero in every interval has zero area. Neutral review explained that
the lowest height on each interval is zero. Every lower estimate of the area
is therefore zero, and integrability gives the result. The fresh retry
noticed the zeros but used only that its lower estimates were at least zero.
Its final inequality cut off. Lean checked both single-point examples and
the general lower-sum argument.

\lead{P02. A missing multiplication.}
The fresh retry placed \(1/n\) at the square indices \(n=1,4,9,\ldots\), with
zeros between them, and showed the series has a finite sum. It stopped just
before checking the product \(n(1/n)\). Lean verified that this product is
one at every square, so it cannot tend to zero. The original rubric credits
that routine calculation. Requiring the last line in print would change
fresh retry from 1/4 repairs to 0/4. Flagged review used height \(1/k\) at
index \(k^2\), whose sum diverges. The audit checked this error as well.

\section{Output budget and interpretation}
\begin{center}
\small\renewcommand{\arraystretch}{1.06}
\begin{tabular}{@{}lrrrr@{}}\toprule
Condition & Replies & Input tokens & Output tokens & Cutoffs\\ \midrule
\costrows
\bottomrule\end{tabular}
\end{center}

On the same four initial errors, A, B and C each produced 8,192 output tokens.
Their input totals were 8,970, 8,918 and 666, respectively. LM Studio recorded
zero API errors. Nine of the 22 cutoffs already contained a valid proof.

Only four problems were eligible for all three follow-ups. One grade changes
a repair rate by 25 percentage points. These are descriptive results for a
handpicked set with one generation per condition per eligible problem.
Repeated draws, independent grading and a larger sample would be needed
to estimate how reliably the pattern carries over.

\clearpage
\appendix
\section{Problem outcomes and replication}
\begin{center}
\small\renewcommand{\arraystretch}{1.07}
\begin{tabular}{@{}llcccc@{}}\toprule
ID & Problem & First & Neutral & Flagged & Fresh\\ \midrule
\itemrows
\bottomrule\end{tabular}
\end{center}
\noindent\small An asterisk marks a reply that reached the 2,048-token cap.
A dash marks a condition outside the allocation. A gap denotes an incorrect
or incomplete proof. P03's first proof was a close cutoff decision. It had
already bounded the two sides of the crossing case before its restatement
ended. The per-answer audit explains the decision.\normalsize

\lead{Model and prompts.}
The local model key was \code{qwen/qwen3.8-27b}, quantized as \code{Q4_K_M}.
The context length was 8,192 tokens, temperature \(0.7\), top-\(p\) \(0.9\),
maximum output 2,048 tokens and requested reasoning effort \code{none}.
The archive includes the model-file hashes, observed load settings and all
request seeds. Each request contained exactly one system message and one
user message. The user instructions were as follows.

\begin{description}\small\setlength{\itemsep}{1pt}
\item[First and fresh.] \quoted{Give a complete proof of the following problem.}
\item[Neutral.] \quoted{Check your previous proof step by step. If you find an error or a gap,
write a complete corrected proof. If the proof is valid, retain it.}
\item[Flagged.] \quoted{This proof has been marked incorrect. Identify any errors or gaps and
write a complete corrected proof.}
\end{description}

\lead{Research materials.}
The public repository contains the ten problems, exact prompts, all 28
responses, blinded grading displays, reasons for each decision, token counts
and the pinned Lean project. The verification script reconstructs each request
and reproduces the tables without calling Qwen. It also checks the saved
review grades against the hashes recorded before unblinding.
\href{https://github.com/vivaanp0011/qwen-proof-self-review}{github.com/vivaanp0011/qwen-proof-self-review}.

\lead{Related work.}
Self-Refine studied iterative feedback and revision across generation tasks
\cite{madaan}. Huang et al.\ tested whether models could correct their own
reasoning \cite{huang}. The fresh-retry condition here asks how much of any
improvement might also occur from another attempt on the same question.

\begin{thebibliography}{2}
\small\raggedright
\bibitem{madaan} A. Madaan et al.\ \emph{Self-Refine: Iterative Refinement with Self-Feedback}.
NeurIPS, 2023. \href{https://arxiv.org/abs/2303.17651}{arXiv:2303.17651}.
\bibitem{huang} J. Huang et al.\ \emph{Large Language Models Cannot Self-Correct Reasoning Yet}.
ICLR, 2024. \href{https://arxiv.org/abs/2310.01798}{arXiv:2310.01798}.
\end{thebibliography}
\end{document}
